% !TITLE 蒿里行
% !AUTHOR 曹操
% !DYNASTY 魏晋
% !GENRE 五言古诗
% !ORDER 1
% !SOURCE 《乐府诗集》卷二十七，相和歌辞·相和曲
% !EDITION 《先秦汉魏晋南北朝诗》，逯钦立辑校，中华书局
% !EDITION_STATUS pending

\begin{poem}
关东有义士，兴兵讨群凶\textsuperscript{1}。
初期会盟津\textsuperscript{2}，乃心在咸阳\textsuperscript{3}。
军合力不齐，踌躇而雁行\textsuperscript{4}。
势利使人争，嗣还自相戕\textsuperscript{5}。
淮南弟称号\textsuperscript{6}，刻玺于北方\textsuperscript{7}。
铠甲生虮虱\textsuperscript{8}，万姓\textsuperscript{9}以死亡。
白骨露于野，千里无鸡鸣。
生民百遗一\textsuperscript{10}，念之断人肠。
\end{poem}

\begin{annotations}
\notetext{1}{义士：指关东起兵讨伐董卓的州郡长官。群凶：指董卓及其党羽。}
\notetext{2}{初期：起初约定。盟津：即孟津，黄河古渡口，在今河南孟津东北。周武王伐纣时曾在此会合诸侯，诗中以盟津代指讨董联军的会合地。}
\notetext{3}{乃心在咸阳：心向汉室。咸阳是秦的都城，这里借指长安，当时汉献帝被董卓挟持在长安。乃心，其心。}
\notetext{4}{踌躇：徘徊不进。雁行：雁飞时排成的行列，形容诸侯军列阵观望，谁也不肯先出兵。}
\notetext{5}{嗣还自相戕：随即自相残杀。嗣还（xuán），随即。戕（qiāng），杀害。}
\notetext{6}{淮南弟称号：指袁术在淮南称帝。袁术是袁绍的从弟，建安二年（197）在寿春称帝。}
\notetext{7}{刻玺于北方：指袁绍在北方刻金玺，图谋废立。玺（xǐ），皇帝的印章。}
\notetext{8}{虮虱：虮（jǐ），虱卵；虱，寄生在人畜身上的小虫。铠甲久不解，故生虮虱。}
\notetext{9}{万姓：百姓。}
\notetext{10}{生民百遗一：百姓百人中只活下一个。遗，留下。}
\end{annotations}

\begin{appreciation}
《蒿里行》是曹操借乐府旧题写时事的名篇之一。蒿里相传是死者魂魄聚居的地方，蒿里之歌本是出殡时唱的挽歌。曹操用这首曲子纪念乱世下的亡魂。

诗从初平元年（190）的关东联军写起。董卓挟着汉献帝西迁长安，关东州郡起兵讨伐，忠臣、义士、野心家们纷纷登场。“关东有义士，兴兵讨群凶”，讲的就是这段历史。曹操自己也在义士之列，起兵后与董卓部将徐荣战于荥阳，几乎全军覆没。“初期会盟津，乃心在咸阳”，联军起初约定在盟津会师，众心都向着长安。盟津是周武王伐纣时会合诸侯的渡口，咸阳借指长安，这两个地名点名了联军打出的正义旗号。

转折在第五句。“军合力不齐，踌躇而雁行”，兵合到一处，心却不齐，谁也不肯先动，像一排大雁停在半空，只摆队形，不往前飞。“势利使人争，嗣还自相戕”，权势和利益一摆上桌，诸侯先自己打了起来。“嗣还”是随即的意思，翻脸翻得这样快，这两个字里有曹操的冷眼。“淮南弟称号，刻玺于北方”，袁术在淮南称帝，袁绍在北方刻金玺、图谋另立天子。讨董的旗子还挂着，旗底下的人已经各自称王称霸，开篇那声“义士”，到这里成了一个反讽。

前八句写诸侯，后八句写苍生。“铠甲生虮虱，万姓以死亡”，仗打得太久，将士的铠甲不解，里面生了虮子和虱子，百姓一批批死去。“白骨露于野，千里无鸡鸣”，荒野上堆着无人收殓的白骨，千里之内听不见一声鸡叫。曹操后来在《军谯令》里回忆故乡，说走了一整天见不到一个熟人，可与这两句对读。“生民百遗一，念之断人肠”，一百个百姓里活不下一个，想到这里，肝肠寸断。

全诗十六句，目光从诸侯落到士兵，从士兵落到百姓，最后一句落回自己身上。《十五从军征》借一个老兵的回家，写战争留给一个家的废墟；这一首铺开的是整个天下。曹操没有出来评论，他只把事实一件件排开，让事实自己说话。钟惺《古诗归》评《薤露行》《蒿里行》为“汉末实录，真诗史也”。

一支本来送给死人的歌，唱出了活人怎样死去。
\end{appreciation}
